\section{Composing Concurrent Commands}
\label{sec:method}

% VALUE: the overall scheme in one paragraph: the composer takes an existing design's per-operator instruction streams, its column map, its dependency graph, and per-operator time estimates, and emits composed commands that run the same kernels, weights, and configuration with bit-identical outputs; it works at the runtime seam and changes no kernel, configuration, compiler, driver, or firmware.
We build a composer that takes an existing design's per-operator instruction streams, its map of operators to columns, the dependency graph of its host program, and an estimate of each operator's device time.
From these, it emits a sequence of composed commands that runs the same kernels, weights, and configuration, and produces bit-identical outputs.
The composer replaces only the host program's dispatch of commands and changes no kernel, configuration, compiler, driver, or firmware.
It makes three kinds of decisions in turn: which operators share a command (Section~\ref{subsec:plan}), the CP's instruction order within each command (Section~\ref{subsec:order}), and where each operator's buffers live (Sections~\ref{subsec:arenas} and~\ref{subsec:rings}).
Each composed command passes a static legality check before any dispatch (Section~\ref{subsec:compose}), and an event model of the CP, the groups, and the host prices a candidate schedule (Section~\ref{subsec:pricing}).
% Implementation: amd-npu a0ee2edf44e0 (tag composer-v1.1), prefill-overlap/{txn.py (111 lines), merged.py (616), component_model.py (180)}; out of tree, Python, on the engine's hardware context.

\subsection{Composing Instruction Streams}
\label{subsec:compose}

% VALUE: a composed command is an interleaving of operators' segments that keeps each operator's own order; it is legal when what the operators touch on the device is disjoint, which is checked statically from the instructions themselves.
\textbf{Segments.}
The composer splits each operator's instruction stream, already reduced to its own columns, into segments as defined in Section~\ref{subsec:formulation}.
A composed command concatenates segments of several operators in a chosen order, keeping each operator's segments in their own order and every instruction exactly once.
Under these rules, an interleaving changes only when each operator's work is issued and awaited, never what the work is.
% txn.segments / txn.compose (amd-npu a0ee2ed prefill-overlap/txn.py): a segment = a run of syncs (0x80) plus the non-sync ops after it; compose asserts per-stream segment order, every op kept once, equal headers.

\textbf{Legality.}
Two operators may interleave in one command only if they cannot disturb each other's state on the device.
The composer checks this from the instructions themselves.
Every register write of an operator must address a column of its own group, and the writes may touch only shim BDs and DMA queues, the operator's own cores' locks, and runtime parameters in its own cores' memory, never core programs, core resets, or memory tiles.
Interleaved operators must patch disjoint argument slots or, as Section~\ref{subsec:arenas} describes, disjoint byte ranges of one slot's buffer.
No operator may write a buffer that another operator in the same command reads or writes.
Operators on the same group, such as several GEMMs, are composed in sequence rather than interleaved, so their writes to that group's registers keep the order of separate commands.
% Each check goes red on a defeated input: swapped groups, an unstripped stream, a sync before an issue op, a patch onto another group's arg, a duplicated column, a reordered or dropped segment (lab-prefill-overlap (3), (4)).

\subsection{Argument Arenas}
\label{subsec:arenas}

% VALUE: existing designs bind one buffer per argument slot, so a command holds only one GEMM; retargeting a patch's byte offset (never its slot) lets any number of GEMMs share the GEMM slots through arenas, with a bound that the planner respects.
The configuration's kernel declares a fixed set of argument slots, and the existing design gives each operator kind its own: three slots for a GEMM's input, weights, and output, and three for each attention group.
With one buffer per slot, a command that carries attention can hold only one GEMM.
An address patch names a slot and a byte offset within that slot's buffer.
We find that retargeting a patch to another slot faults in the IOMMU, while changing its byte offset is safe and keeps outputs bit-identical.
The composer therefore allocates the GEMM inputs, weights, and outputs as sub-buffers of three parent buffers, which we call \emph{arenas}, bound to the GEMM slots, and rewrites only each GEMM's patch offsets.
Any number of GEMMs then share three slots.
The offset field bounds an arena to $2^{31}$ bytes, so the composer splits the weights into arenas of whole layers below that bound, and the GEMMs of one command must share a weight arena.
% Slots: nargs = 3 + 3 x len(cfg.attn) = 9 (mlir-air a3e26261 engine.py:107). Patch: arg index at word 8, byte offset at word 10 of opcode 0x81.
% Retarget checks (lab-prefill-overlap (8), amd-npu 61c64f9 prefill-overlap/retarget/): slot moves complete with wrong output and AMD-Vi IO_PAGE_FAULT; offset moves inside a parent BO are bit-identical; two GEMMs at distinct offsets in shared arenas are bit-identical.
% Arena bound: merged.Weights (a0ee2ed); Llama-3.2-1B: one weight arena; Qwen3-4B: two (2145 + 195 MB) (lab-prefill-overlap (24), (26)).

\subsection{Planning Commands}
\label{subsec:plan}

% VALUE: the planner list-schedules the dependency graph chunk by chunk rather than stage by stage, so each attention operator meets independent GEMMs of earlier chunks, and it sizes each command so the GEMM work matches the attention work.
The existing host program runs each layer stage by stage, so every output projection waits behind the layer's entire attention, and no GEMM is ready to run beside attention.
The composer instead list-schedules the dependency graph chunk by chunk.
At each step, it takes the ready attention operator of the lowest layer and chunk and adds ready GEMMs, in the same order, until their estimated work reaches that of the attention operator.
The GEMMs then run in sequence on columns 1--6 while attention runs on columns 0 and 7.
A step without a ready attention operator packs up to a fixed number of ready GEMMs into one command, following standard practice for amortizing dispatch cost.
Operators separated by host work never share a command, and a GEMM whose weights lie in another arena waits for a later command.
% merged.plan (a0ee2ed): deps qkv(L,c) <- down(L-1,c); attn(L,c) <- qkv(L,0..c); o <- attn; gate, up <- o; down <- gate, up. Optional cap K on GEMMs riding with attention; B GEMMs per command without attention.

\subsection{Instruction Order}
\label{subsec:order}

\begin{figure}[t]
    \centering
    % VALUE: in one command, the in-turn order makes the CP wait on attention before it issues the second GEMM, while the completion order issues each GEMM as soon as the previous one completes and keeps both attention columns fed.
    \figph{Schematic: CP instruction order and column activity for one command (attention on columns 0 and 7, three GEMMs on columns 1--6), in turn vs.\ in completion order.}
    \caption{CP instruction order and the resulting activity of the column groups in one composed command: (a) the operators in turn; (b) synchronizations in predicted completion order.}
    \label{fig:order}
\end{figure}

% VALUE: within a command, the CP opens every group's work first and then places each synchronization at its predicted completion, followed by the issue it gates, so no synchronization holds back work for another group; the order is fixed before dispatch from the time estimates.
Within a command, the order of synchronizations decides how long each group waits for its next issue (Figure~\ref{fig:order}).
In turn, the CP would issue the first GEMM, then walk the attention stream to its end, and only then issue the second GEMM, so the GEMM group idles behind attention's synchronizations.
The composer instead issues the opening segments of both attention columns and of the first GEMM, and then places every synchronization in order of its \emph{predicted completion}, each followed by the issue it gates.
A GEMM's synchronizations are predicted to complete after the sum of the estimates of the GEMMs before it, and attention's mid-stream synchronizations at fixed fractions of its estimate.
Attention's two columns are treated as two groups, so their segments alternate, which removes the serialization that Section~\ref{subsec:semantics} measured inside the attention operator.
The order is static: the estimates fix it when the command is built, and the device executes it without host involvement.
When issue costs are negligible, the estimates are exact, and the host work is hidden, the CP never blocks on a later-completing operator while an earlier one's successor waits, so the order matches the list schedule of the groups without a CP.
% merged.cp_order (a0ee2ed): attention mid-sync fractions MID = (0.25, 0.5); measured best among (0.1,0.2), (0.25,0.5), (0.4,0.7), (0.6,0.85), (0.9,0.95); the in-turn order costs +26% on the same 38 commands (lab-prefill-overlap (16)).
% TODO(proof): the exactness statement is argued, not proven (ch3-formulation § Lower bound and heuristic).

\subsection{Buffer Rings}
\label{subsec:rings}

% VALUE: with commands in flight and chunks interleaved, each buffer role needs several slots; the composer derives each ring's depth from the plan's live intervals and checks that no slot is reused before its reader is done.
Once operators of several chunks are in flight, a GEMM's input and output buffers cannot be reused per stage as in the existing host program.
For each buffer role, the composer derives a ring of slots in the arenas from the plan: each use of a buffer lives from its producer's write to its last reader's completion, and interval colouring assigns slots so that no two overlapping uses share one.
A separate check asserts, over the whole plan, that no slot is reused before its previous reader has completed.
The depths follow from the plan and are not set by hand.
For example, every QKV input of a layer is written before the layer's first QKV command completes, so that ring needs one slot per chunk.
% merged.colour / assert_no_hazard / Layout (a0ee2ed); the hazard check goes red on an overlapping slot. Rings at 4k: shared 2048-wide input 32, 8192-wide input 4, outputs 4-18 (lab-prefill-overlap (9)).

\subsection{Host Boundaries}
\label{subsec:host}

% VALUE: host work between commands forces boundaries; keeping D commands in flight hides the host and the boundary cost but delays which operators are ready, so the depth and the packing trade overlap against readiness.
The existing host program computes normalizations, rotary embeddings, and residuals on the CPU, so a dependence that crosses the host forces a command boundary.
The composer keeps $D$ commands in flight: it queues the next command before it waits on the oldest, and it post-processes a command's outputs while later commands run.
A command can therefore carry only operators whose inputs come from commands at least $D$ positions earlier.
A deeper queue hides more of the boundary cost and of the host work, but it leaves fewer operators ready to share a command with attention.
The composer models this delay directly, so the queue depth and the packing of GEMMs are chosen per prompt length (Section~\ref{sec:eval}).
% merged.plan lag = D; Layout: an output slot is claimed when its command is queued, D - 1.6 commands before it runs.

\subsection{Pricing a Schedule}
\label{subsec:pricing}

\begin{table}[t]
    \centering
    \caption{Parameters of the Event Model, Llama-3.2-1B, 2{,}048 Tokens}
    \label{tab:params}
    % VALUE: every term of the model is measured, on the device or from the host's timestamps, and none is a free scale factor.
    \begin{threeparttable}
    \begin{tabular*}{\linewidth}{@{\extracolsep{\fill}}>{\raggedright\arraybackslash}p{0.42\linewidth}lr}
        \toprule
        Term & Source & Value \\
        \midrule
        % SOURCE: generated by figs/make_figs.py from figs/data/component-model.json; do not edit.
\textit{Device} & \\
GEMM busy (QKV, O, gate, up, down) & 0.14 / 0.10 / 0.51 / 0.40 / 0.43\ms \\
Attention busy at chunk $c$ & $0.062 + 0.072\,c$\ms \\
Gap per additional GEMM $\gamma$ & 35\us \\
Overlap cost $\beta$ & 0.12 \\
\midrule
\textit{Host} & \\
Boundary cost $\epsilon$ (attention, GEMM only) & 50\us / 57\us \\
Issue cost (attention, GEMM only) & 39\us / 9\us \\
Post cost (QKV, attention, O) & 0.25 / 0.06 / 0.07\ms \\
Post cost (gate, up, down) & 0.03 / 0.20 / 0.07\ms \\
Head and tail & 22.9\ms \\
\bottomrule

        \bottomrule
    \end{tabular*}
    \begin{tablenotes}\footnotesize
    \item Device terms from each command replayed alone; host terms from per-command timestamps of three schedules at 2{,}048 tokens.
    \end{tablenotes}
    \end{threeparttable}
\end{table}

% VALUE: to choose among schedules without running each one, the composer prices a plan with an event simulation whose every term is measured; it separates device time, host exposure, and the fixed head and tail.
To choose among plans without running each on the device, the composer prices a plan with an event simulation of the device, the CP, and the host (Table~\ref{tab:params}).
On the device, command $k$ ends at
\begin{equation}
\label{eq:cmd}
e_k = \max\bigl(e_{k-1} + \epsilon,\; q_k + F\bigr) + b_k,
\end{equation}
where $q_k$ is the time the host queues it, so a queued command pays the boundary cost $\epsilon$ and a command that arrives at an idle device pays the dispatch floor $F$.
A command's busy time $b_k$ is the sum of its GEMMs' busy times plus a gap $\gamma$ per additional GEMM, and a command that also carries attention takes $\max(a, g) + \beta \min(a, g)$ for attention work $a$ and GEMM work $g$.
Attention's busy time grows linearly with the chunk index, as the number of earlier key blocks does.
The host pays an issue cost per command and a post-processing cost per operator, and waits on the oldest command once $D$ are in flight.
The TTFT is the simulated span plus a fixed head and tail, namely the embedding, the first normalization, and the LM head.
Because each term is measured separately, an error in one term cannot hide behind another, and the model reports the device time, the host's exposed time, and the head and tail separately.
% component_model.py (amd-npu a0ee2ed prefill-overlap/component_model.py; parameters prefill-overlap/v1/component_model.json, first committed at 6735c61). Fit arms 2k depth 2/4/16 (no packing); held-out arms are the evaluation's.
